\begin{table*}[htbp]
\centering
\setlength{\tabcolsep}{8pt}
\renewcommand\arraystretch{1.0}
\begin{tabular}{c|c|ccc|ccc|ccc}
\toprule
\multirow{2}{*}{\textbf{Setting}}
 & \multirow{2}{*}{\textbf{Method}}
 & \multicolumn{3}{c|}{\textbf{Crime}}
 & \multicolumn{3}{c|}{\textbf{Crash}}
 & \multicolumn{3}{c}{\textbf{Nightlight}} \\ \cmidrule(lr){3-5}\cmidrule(lr){6-8}\cmidrule(lr){9-11}
 &  & MAE & RMSE & PCC & MAE & RMSE & PCC & MAE & RMSE & PCC \\
\midrule
\multirow{9}{*}{\textbf{NYC$\rightarrow$ CHI}}
 & TS2Vec             & \underline{873.36}      & 1192.75      & 0.49   & 334.72      & 478.50      & 0.45   & 18.18      & 23.83      & 0.41 \\
 & PatchTST         & 1608.18      & 2062.95      & 0.45   & 585.96      & 735.64      & 0.50   & 21.84      & 33.49      & 0.25 \\
 & TSLANet                & 1609.46      & 2064.82      & 0.17   & 698.64      & 864.23      & 0.31   & 19.33      & 26.72      & -0.31\\
 & TimesURL                & 891.21      & 1217.53      & 0.51   & 345.86      & 491.47      & 0.45   & 18.53      & 24.91      & 0.31 \\ \cmidrule(lr){2-11}
 & TransGTR             & 893.54      & \underline{1135.31}      & 0.48   & 329.64      & 473.65      & 0.37   & 17.34      & 23.25      & 0.43\\
 & CrossST                 & 960.13      & 1254.25      & 0.49   & 322.90      & 405.74      & 0.55   & 17.36      & 23.37      & 0.47\\
 & CoRE                 & 1631.15      & 2082.45      & 0.17   & 569.55      & 729.45      & 0.26   & \underline{16.07}      & \underline{20.20}      & \underline{0.59}\\
  & UrbanPG             & 947.09      & 1200.80      & \underline{0.51}   & \underline{318.97}      & \underline{399.68}      & \underline{0.59}   & 16.63      & 22.24      & 0.53\\ \cmidrule(lr){2-11}
 & MAP-STAR     & \textbf{864.53}  & \textbf{1105.60} & \textbf{0.54} & \textbf{302.97}  & \textbf{396.00}  & \textbf{0.63} & \textbf{15.61}   & \textbf{19.21}   & \textbf{0.71}\\
\midrule
\multirow{9}{*}{\textbf{CHI$\rightarrow$NYC}}
 & TS2Vec             & 372.61      & 524.24      & 0.35   & 81.05      & 114.23      & 0.41   & 29.08      & 48.65      & 0.50\\
 & PatchTST         & 417.68      & 603.17      & -0.17   & 95.15      & 132.91      & 0.02   & 37.77      & 51.98      & 0.42\\
 & TSLANet                & 384.10      & 537.64      & 0.33   & 87.23      & 121.09      & 0.31   & 30.67      & 49.17      & 0.47\\
 & TimesURL                & 376.06      & 536.21      & 0.29   & 79.50      & 113.65      & 0.41   & 29.39      & 47.03      & 0.48\\ \cmidrule(lr){2-11}
 & TransGTR             & 372.92      & 522.92      & 0.30   & 78.43      & 109.02      & 0.48   & 26.93      & 43.67      & 0.56\\
 & CrossST                 & 369.25      & 494.47      & 0.46   & \underline{75.47}      & 106.35      & 0.48   & 26.63      & 42.88      & 0.61\\
 & CoRE                 & 451.26      & 669.16      & 0.01   & 85.31      & 119.85      & 0.26   & 28.50      & 48.02      & 0.44\\
  & UrbanPG             & \underline{357.92}      & \underline{483.76}      & \underline{0.49}   & 76.17      & \underline{104.28}      & \underline{0.55}   & \underline{25.45}      & \underline{42.76}      & \underline{0.63}\\ \cmidrule(lr){2-11}
 & MAP-STAR     & \textbf{340.79}  & \textbf{476.99}  &  \textbf{0.51}  & \textbf{73.49}  & \textbf{101.06} & \textbf{0.58} & \textbf{25.04}   & \textbf{41.06}   & \textbf{0.66}\\
\bottomrule
\end{tabular}
\caption{
% Performance comparison with baselines in NYC2CHI and CHI2NYC cross-city scenarios on urban computing tasks.
Cross-city tasks comparison results. “A$\rightarrow$B” denotes that the models are trained on city A and evaluated on city B.
}
\label{tab:cross_city}
\end{table*}

\section{Experiment}

In this section, we aim to answer the following research questions through the experimental evaluations:

\begin{itemize}
    \item RQ1: How does MAP-STAR perform compared to state-of-the-art region representation learning methods in different cross-city downstream tasks?
    \item RQ2: How do components in MAP-STAR affect overall performance in transferable representation learning?
    % RQ3: 相比于传统的在全图上进行区域表示学习的方法，MAP-STAR基于异质子图采样的表示学习能否学到充足的区域结构和语义信息？
    % RQ4: 在地理／文化背景多样的条件下，MAP-STAR 在更广泛下游任务中的迁移性如何？
    % RQ5: 不同的子图采样策略和子图采样规模对模型效果会产生怎样的影响？
    \item RQ3: Compared with traditional full-graph-based representation learning methods, can MAP-STAR’s heterogeneous subgraph sampling still capture sufficient regional information in structure and semantics?
    \item RQ4: How transferable is MAP-STAR across cities with substantially diverse geographic and cultural contexts in a broad range of downstream tasks?
    \item RQ5: How do different subgraph sampling strategies and subgraph scales affect the model’s performance?
    \item RQ6: Do models pre-trained on different cities yield consistent representation distributions and semantic clustering in the target domain?
\end{itemize}



\subsection{Experimental Setup}
\subsubsection{Datasets.} 

We conduct experiments using taxi data from New York City (NYC) and Chicago (CHI) as pre-training sources. Regional inflow and OD statistics are aggregated hourly based on region boundaries from the NYC TLC and CHI Data Portal Community Areas. We then collect POI data from OpenStreetMap (OSM) and construct heterogeneous subgraphs $\mathcal{G}_i^{\text{sub}}$ by integrating region boundaries, taxi flows and POI distributions. For downstream tasks, we collect crime, crash, and business datasets from NYC Open Data and CHI Data Portal, along with nighttime light data from NPP-VIIRS-like NTL data~\cite{DVN/YGIVCD_2020}, all aligned with the pre-trainingng period. These support four evaluation tasks that assess the quality and transferability of the representations of the learned regions.

\subsubsection{Baselines \& Evaluating Metrics.}

We compare MAP-STAR against eleven baselines from three categories under two experimental settings. In the single-city setting, we evaluate MAP-STAR alongside traditional graph-based region representation learning methods, including HREP~\cite{HREP}, ReCP~\cite{ReCP}, and ST-ReP~\cite{ST-ReP}. In the cross-city setting, we benchmark against two additional groups: sequence-based methods that encode traffic flow sequences as region representations, including TS2Vec\cite{TS2Vec}, PatchTST~\cite{PatchTST}, TSLANet~\cite{TSLANet}, and TimesURL~\cite{TimesURL}; and region transfer learning methods that aim to transfer spatio-temporal knowledge from a source city to a target city, including TransGTR~\cite{TransGTR}, CrossST~\cite{CrossST}, CoRE~\cite{CoRE} and UrbanPG~\cite{UrbanPG}. For all regression-based tasks, we report MAE, RMSE, and PCC, where lower MAE and RMSE indicate better performance and higher PCC reflects stronger predictive alignment.

\subsubsection{Implementations.}

We implement MAP-STAR using the Adam optimizer with a fixed learning rate of 0.001 and a representation dimension of 128. During pretraining, we segment the input into patch sequences with a patch length of 2016 and a sliding window of 168. For subgraph construction, we sample second-order neighbors for each region node, with both first-order and second-order neighbor counts set to 5 to ensure a balanced local topology. An early stopping strategy is applied with a patience of 30 during pretraining and 50 during downstream evaluation. Detailed hyperparameter settings are provided in the Appendix.

\begin{figure}[!htbp]
   \centering
    \includegraphics[width=0.50\textwidth]{fig/ablation.png} % Reduce the figure size so that it is slightly narrower than the column.
    \caption{Ablation Study of the proposed MAP-STAR.}
    \label{fig:ablation}
\end{figure}

\subsection{Performance Comparison (RQ1)}

In this section, we conduct comparative experiments in cross-city settings against two categories of baselines. From the result show in Table~\ref{tab:cross_city}, we observe that:
\begin{itemize}

    \item Traditional time-series representation models consistently lag behind existing region transfer-learning approaches. The main reason is that human mobility are highly city-specific. Representations trained only on a source city struggle to generalise to a different urban context. A cross-task comparison further reveals that such baselines fare relatively better on the crash prediction task than on crime or nightlight. Traffic accidents are driven largely by mobility patterns that classic temporal models can still capture. Whereas crime and nocturnal luminosity are more closely tied to region functional properties like POI distributions. This contrast underscores the necessity of modelling heterogeneous inter-regional graph structures and their interactions when modeling transferable region representations.
    
    \item Graph-based transfer-learning methods outperform traditional time-series representation models on most downstream tasks. To compensate for the latter’s inability to capture interregional mobility, TransGTR constructs an adjacency graph of flow similarities, allowing the model to learn part of the interaction patterns between regions. CrossST goes a step further by explicitly modelling recurrent traffic-flow regularities along both the temporal and spatial dimensions. These designs alleviate the limitations of conventional time-series approaches in sequence-pattern transfer and graph-structure modelling. However, relying solely on traffic flow data and ignoring structural heterogeneity between regions still constrains the expressive power of the representations.

    \item The proposed MAP-STAR consistently outperforms all baselines, demonstrating its ability to model multi-view inter-regional interactions and learn representations rich in diverse regional semantics. Unlike CrossST, which applies a fixed spatial encoding across cities' different spatial dimensions, MAP-STAR employs a heterogeneous subgraph sampler to extract identical subgraph structures in both cities. This stabilises the captured multi-view patterns and further boosts the model’s overall performance.
\end{itemize}

\begin{table}[htbp]
\centering
\setlength{\tabcolsep}{2.75pt}
\renewcommand\arraystretch{1.31}
\begin{tabular}{c|cc|cc|cc}
\toprule
Task & \multicolumn{2}{c|}{\textbf{Crime}} & \multicolumn{2}{c|}{\textbf{Crash}} & \multicolumn{2}{c}{\textbf{Nightlight}} \\
\midrule
Metrics & MAE & RMSE & MAE & RMSE & MAE & RMSE \\
\midrule
HREP & \underline{345.37} & \underline{476.07} & 78.04 & 108.44 & 29.86 & 47.18 \\
ReCP  & 349.79 & 506.87 & \underline{69.45} & \underline{95.57} & \textbf{25.52} & \underline{44.62} \\
ST-ReP  & 385.65 & 558.92 & 85.05 & 117.77 & 28.93 & 46.89 \\
MAP-STAR  & \textbf{334.92} & \textbf{461.22} & \textbf{64.14} & \textbf{91.10} & \underline{25.65} & \textbf{42.25} \\
\bottomrule
\end{tabular}
\caption{Performance comparison in single-city scenario.}
\label{tab:single_city}
\end{table}

\subsection{Ablation Study (RQ2)}

To assess the contribution of each component in MAP-STAR, we conduct an ablation study under the CHI$\rightarrow$NYC transfer setting. We compare three reduced variants in terms of RMSE: (1) w/o FFT, which omits Fourier transform and inverse transform during temporal encoding; (2) w/o MRGCN, which restricts message passing to geographic neighbors only and removes the multi-relation message-passing (MRMP) mechanism; and (3) w/o MVPBA, which feeds the representations produced by MRGCN directly into the reconstruction decoder, bypassing the multi-view predictive bidirectional alignment (MVPBA) module.

As illustrated in Figure \ref{fig:ablation}, MAP-STAR consistently outperforms all ablated variants, demonstrating the necessity of the full architecture. Removing MRMP prevents nodes from aggregating messages from semantically similar or mobility-connected but geographically distant nodes, thereby degrading representation quality. Eliminating MVPBA weakens cross-domain alignment and regularization, which lowers performance on cross-city tasks. Moreover, incorporating FFT filtering in the Time Encoder enables the model to capture generic mobility cycles and long-term trends, further enhancing robustness and accuracy in transfer scenarios.


\subsection{Effectiveness for Sub-Graph Modeling (RQ3)}

To examine whether heterogeneous subgraph sampling can rival—or even surpass—conventional full-graph training in a single-city context, we conduct a focused study on the New York City (NYC) dataset. Specifically, we compare MAP-STAR with three state-of-the-art region-representation baselines: HREP\cite{HREP}, ReCP\cite{ReCP}, and ST-ReP\cite{ST-ReP}. The training pipeline for MAP-STAR mirrors the cross-city protocol; the sole alteration is that both the pre-training and fine-tuning phases are executed exclusively on NYC data.

The results in Table~\ref{tab:single_city} reveal a consistent advantage for MAP-STAR.  These gains confirm that sampling multiple small, semantically-complementary subgraphs is sufficient to capture rich regional semantics, even when the global graph is fully accessible.  We attribute the improvements to two factors: (i) heterogeneous views allow the model to focus on neighbourhood, POI similarity and flow-based interactions separately, avoiding the noise of a monolithic graph; and (ii) the pattern-bank adapter re-uses structural motifs shared across views, yielding more compact yet expressive representations.  Collectively, the single-city study demonstrates that MAP-STAR is not only transferable but also a lightweight and generalizable alternative to full-graph encoders for region representation learning.

\begin{figure*}[t]
\centering
\includegraphics[width=0.99\textwidth]{fig/visualization_v3.png} % Reduce the figure size so that it is slightly narrower than the column.
% 
\caption{Representation Visualization Results in both single-city and cross-city settings.}
\label{fig:vr2}
\end{figure*}

\begin{table}[htbp]
\centering
\setlength{\tabcolsep}{2.8pt}
\renewcommand\arraystretch{1.25}
\begin{tabular}{c|ccc|ccc}
\toprule
Setting & \multicolumn{3}{c|}{\textbf{CHI$\rightarrow$NYC}} & \multicolumn{3}{c}{\textbf{CHI$\rightarrow$CD}} \\
\midrule
Metrics & MAE & RMSE & PCC & MAE & RMSE & PCC \\
\midrule
TransGTR & 15.42 & 19.59 & 0.57 & 260.19 & 467.40 & 0.40 \\
CrossST  & 16.39 & 22.71 & 0.52 & 214.99 & 338.54 & 0.68 \\
MAP-STAR & \textbf{14.36} & \textbf{18.05} & \textbf{0.58} & \textbf{196.40} & \textbf{307.03} & \textbf{0.77} \\
\bottomrule
\end{tabular}
\caption{Generalization performance across cities.}
\label{tab:generalization}
\end{table}

\subsection{Generalization in Cross-City Transfer (RQ4)}

To further evaluate the transferability of MAP-STAR under heterogeneous urban conditions and a broader downstream task, we conduct a challenging cross-city generalization study. We consider business POI count prediction as the downstream task, where models are pre-trained on Chicago (CHI) and evaluated on New York City (NYC) and Chengdu (CD). Compared with CHI, Chengdu exhibits substantial differences in urban scale, POI distributions, and mobility patterns, resulting in a more heterogeneous transfer scenario. We report results for MAP-STAR and two strong baselines, TransGTR and CrossST, in Table~\ref{tab:generalization}.

As shown in the results, MAP-STAR consistently outperforms the baselines across all transfer settings, indicating stronger generalization under cross-city distribution shifts. In particular, TransGTR shows a pronounced performance drop in the CHI$\rightarrow$CD setting, suggesting that distillation-based transfer becomes less robust when graph structures and data distributions differ substantially. CrossST yields more stable results yet still falls short of MAP-STAR, particularly under the more challenging CHI$\rightarrow$CD scenario. In contrast, MAP-STAR leverages heterogeneous subgraph sampling and a multi-view structural pattern bank to capture reusable structural patterns rather than city-specific features, thereby learning more stable and transferable representations across diverse urban graphs.

\begin{figure}[!htbp]
   \centering
    \includegraphics[width=0.485\textwidth]{fig/graph_sampler.png} % Reduce the figure size so that it is slightly narrower than the column.
    \caption{Empirical results on different sampling settings.}
    \label{fig:graphsample}
\end{figure}

\subsection{Experiments for Subgraph Sampler (RQ5)}

We conduct an extensive study to examine how different subgraph sampling strategies and subgraph scales affect performance under the CHI$\rightarrow$NYC transfer setting. Specifically, the \textit{OD} strategy samples subgraphs solely based on inter-regional population flows, \textit{POI} relies on similarity in points-of-interest distributions, \textit{Neighbor} considers only geographically adjacent regions, while \textit{Mix} combines all three views into a heterogeneous subgraph.

As shown in Figure \ref{fig:graphsample}(a)–(c), the \textit{Mix} strategy consistently achieves the best performance, indicating that heterogeneous subgraphs enable the model to aggregate complementary relational signals and learn richer region semantics. An exception is observed in the Nightlight task, where POI-based sampling yields the second-best results, suggesting that nighttime luminosity is more strongly associated with intrinsic regional attributes than with mobility patterns or spatial adjacency. This observation further implies that different downstream tasks may favor different relational views. We also analyze the effect of subgraph scale, where $k$ denotes the number of sampled neighbors per hop within a two-hop neighborhood. As illustrated in Figure \ref{fig:graphsample}(d), performance peaks when $k$ lies in the range of 4–5: smaller subgraphs provide insufficient context, whereas larger ones tend to introduce excessive noise. Based on these results, we set the subgraph scale to $k=5$ in all experiments.

\subsection{Representation Visualization (RQ6)}

To further examine the quality of the representations learned by MAP-STAR, we visualize region embeddings obtained under both single-city and cross-city settings. Specifically, MAP-STAR is pre-trained on NYC and CHI, and embeddings are extracted from models trained on NYC alone as well as models transferred from CHI to NYC. These embeddings are then projected into a two-dimensional space using dimensionality reduction techniques for qualitative comparison. To assess whether heterogeneous subgraph sampling enables the model to capture urban semantics from multiple complementary views, we visualize the embeddings colored by four region-level labels that reflect distinct urban properties: most frequent POI category, population flow intensity, latitude, and longitude.

The results are shown in Figure \ref{fig:vr2}. After dimensionality reduction, embeddings associated with different labels form clearly separated clusters, indicating that MAP-STAR effectively encodes functional similarity, mobility dynamics, and geographic proximity into a unified representation space without conflating these distinct semantic dimensions. Moreover, embedding distributions under different training settings exhibit consistent structures on the same target city. For instance, in both the "NYC" and "CHI$\rightarrow$NYC" settings, longitude-colored visualizations reveal a similar outlying cluster corresponding to geographically extreme regions. The preservation of such fine-grained spatial structures demonstrates that MAP-STAR learns transferable region representations that capture stable urban characteristics regardless of the pre-training source city, highlighting its robustness under cross-city distribution shifts.